\subsection{Welfare and sale terms}\label{app:welfare}

This subsection states and proves Propositions~\ref{prop:access}, \ref{prop:payment} and~\ref{prop:reserves}, and
it proves the statements of Section~\ref{sec:welfare} that rest on them.

\subsubsection*{Access to prices}

Hold the incumbent at \(F_1\). Compare the equilibrium of Proposition~\ref{prop:ccc} with the equilibrium of an
economy in which the rival does not see the price. All other primitives stay fixed, and market makers set
competitive prices in both economies. Write \(\mathcal R\) for expected target proceeds and \(\mathcal W\) for expected
acquisition surplus net of preparation costs. The superscript \(1\) marks price access, and \(0\) marks the hidden
price. Acquisition surplus excludes transfers between bidders, shareholders and traders.

\begin{proposition}[Access to prices]\label{prop:access}
Let the hypotheses of Proposition~\ref{prop:ccc} hold. At \(F_1\), both economies have unique full orders. Price
access raises expected target proceeds and acquisition surplus net of preparation costs:
\[
\begin{aligned}
\mathcal R^1-\mathcal R^0&=(1-\rho)\bigl\{\pi(b+\lambda)(t_H-t_0)+(1-\pi)\,b\,(t_L-t_0)\bigr\}>0,\\
\mathcal W^1-\mathcal W^0&=(1-\rho)(b+\pi\lambda)\bigl[B_1(M)-c_H+pF_1(p)\bigr]>0.
\end{aligned}
\]
\end{proposition}

\begin{proof}[Proof of Proposition~\ref{prop:access}]
\emph{The two tuples.} The access tuple is \(\mathcal E_{\rm F}(F_1)\) of the proof of Proposition~\ref{prop:ccc}. The hidden
tuple has full orders, the posteriors \((m,\pi,M)\) at the orders, the rival's belief \(\pi\) at every price
because the rival sees neither the price nor the order, \(\zeta\equiv0\), entry \(\rho\) after every order, the
prices \(P^0(x)=t_0+\rho\,w(\mu(x))\), and truthful bids. Here and below every object is at \(F_1\).

\emph{The conditions of the hidden tuple.} (E1): the full-order likelihoods. (E2): the rival's belief is its
prior. (E3): \(B_1(\pi)\le B_0(\pi)<c_H\), strictly. (E4): \eqref{eq:a2-terminal} with \(e\equiv\rho\). (E5): with
\(e\equiv\rho\), the purchase earns \(\rho\Delta_1(1-M)-k\ge\rho\kappa\Delta_1-k>0\) and the sale
\(\rho\Delta_1m-k>0\), and the wrong-signed orders earn at most \(-k\). (E6): truthful bids. These bounds hold for
every candidate of the hidden economy, so its orders are unique, and the tuple is unique. The access economy is
unique by Proposition~\ref{prop:ccc}.

\emph{The proceeds gain.} Couple the two economies on the class, the value, the cost, the slot identity and the
liquidity order. The orders coincide at every realization, so the trading costs coincide. Access adds high-cost
preparation exactly after the order \(1\), which has probability \(b+\lambda\) in class \(H\) and \(b\) in class
\(L\). With entry, expected proceeds rise from \(t_0\) to \(t_j\). Hence
\[
\mathcal R^1-\mathcal R^0=(1-\rho)\{\pi(b+\lambda)w_H+(1-\pi)\,b\,w_L\}.
\]
Each term is nonnegative, and \(w_H=w_L+\Delta_1\ge\Delta_1>0\), so the gain is positive.

\emph{The surplus gain.} Let \(\mathcal V_0(R)=R\1\{R\ge p\}\) be the allocation value without the rival.
With truthful bids, the target goes to the highest admissible value, so with an admissible rival of value
\(\theta\) the allocation value is \(\max\{\mathcal V_0(R),\theta\}\), and
\begin{equation}\label{eq:a7-welfare}
\max\{\mathcal V_0(R),\theta\}-\mathcal V_0(R)=(\theta-\max\{p,R\})_++p\1\{R<p\}.
\end{equation}
For \(R<p\), both sides equal \(\theta\), since \(\theta\ge p\); for \(R\ge p\), both sides equal \((\theta-R)_+\).
Average over \(R\) and over the class at the posterior \(M\) of the order \(1\): one added preparation adds
\(B_1(M)+pF_1(p)\) to allocation value and costs \(c_H\). Added preparations have probability
\((1-\rho)\Pr(X=1)=(1-\rho)(b+\pi\lambda)\). Hence
\(\mathcal W^1-\mathcal W^0=(1-\rho)(b+\pi\lambda)[B_1(M)-c_H+pF_1(p)]>0\), because \(B_1(M)>c_H\). Transfers between
bidders, shareholders and traders cancel from this sum.
\end{proof}

\paragraph{The example, the transfer and the trade scale.}
At \(r_1\), proceeds are \(\reg{p2_base_r1_proceeds}\) with access and \(\reg{p2_base_hidden_proceeds}\) without, a gain of
\(\reg{p2_access_proceeds_gain}\); net surplus is \(\reg{p2_base_r1_surplus}\) and \(\reg{p2_base_hidden_surplus}\), a gain of
\(\reg{p2_access_surplus_gain}\). If the execution cost is a resource cost, both economies lose \(\lambda k\), and the
difference is unchanged. In the access economy the informed trader arrives with probability \(\lambda\); class \(H\) buys
with entry one after a purchase, and class \(L\) sells with entry \(\rho\) after a sale. By the payoffs of
Lemma~\ref{lem:trading} before the cost \(k\), its expected gross profit is
\[
G=\lambda\Delta_1\{\pi(1-M)+(1-\pi)\rho m\}=\reg{p2_access_gross_transfer}.
\]
Competitive prices give \(\E[V_T-P(X)\mid X]=0\), so \(\E[X(V_T-P(X))]=0\). The informed and the liquidity orders
partition this expectation, so liquidity traders lose \(G\) in expected gross terms. Measure one order unit in shares per share of float, and
keep the same signal experiment. The economy without an informed trader has the
posterior \(\pi\) after every order and entry \(\rho\), so its proceeds equal \(\mathcal R^0\). If initial shareholders
bear the whole transfer, their total gain over that economy is \(\mathcal R^1-\mathcal R^0\) minus the order unit
times \(G\). It is positive exactly when the order unit is below \((\mathcal R^1-\mathcal R^0)/G=\reg{p2_access_break_even}\). A deterministic external dividend raises \(V_T\)
and every price by its amount. It leaves \(V_T-P\), every order payoff, the rival's information and entry
unchanged, and it enters neither proceeds of the sale nor allocation surplus. With research at the cap,
Appendix~\ref{app:research} gives the same gains.

\paragraph{Access outside the recruitment region.}
Suppose instead \(B_1(m)<c_H<B_1(\pi)\) and \(k<\rho\kappa\Delta_1\). Without the price, the rival's belief is
\(\pi\), so every rival prepares, and entry is one after every order. With the price, the posteriors are
\((m,\pi,M)\), and entry is \((\rho,1,1)\). In both economies the bound on every profile forces full orders, so the
orders, the posteriors and the trading costs coincide, and each tuple satisfies (E1) to (E6) as above. Access removes
high-cost preparation exactly after a sale, which has probability \(b\) in class \(H\) and \(b+\lambda\) in class
\(L\). Proceeds change by \(-(1-\rho)\{\pi b\,w_H+(1-\pi)(b+\lambda)\,w_L\}<0\), and surplus by
\(-(1-\rho)[b+(1-\pi)\lambda]\,[B_1(m)-c_H+pF_1(p)]\). Keep the strong incumbent of the example and set
\(c_H=\reg{access_out_cH}\). Then \(B_1(m)-c_H=\reg{p2_outside_seen_cost_sign_minus}\) and
\(B_1(\pi)-c_H=\reg{p2_outside_seen_cost_sign_zero}\), and access changes proceeds by \(\reg{p2_outside_proceeds_gain}\)
and surplus by \(\reg{p2_outside_surplus_gain}\).

\subsubsection*{The payment rule}

Consider a verifiable-value institution with zero reserve. After diligence, values are verifiable, and a sale to the next bidder at its value is an enforceable fallback. The seller and the winner divide the surplus above the
fallback by Nash bargaining, with seller weight \(\omega\in[0,1)\). Without a rival, the seller receives
\(\omega R\).

\begin{proposition}[The payment rule]\label{prop:payment}
Let the rival's values be the point masses \(\ell\) and \(h\), and let the incumbent distribution have support
\([0,\bar r]\) with \(0<\ell<\bar r<h\). The transfer, the rival's expected profit and the spread are
\[
\begin{aligned}
T_\omega(R,\theta)&=(1-\omega)\min\{R,\theta\}+\omega\max\{R,\theta\},\\
g_{\theta,\omega}&=(1-\omega)\,\E[(\theta-R)_+],\\
\Delta_\omega&=\omega(h-\ell)+(1-2\omega)\,\E[(R-\ell)_+].
\end{aligned}
\]
A stronger incumbent weakly lowers the rival's profit. Its effect on the spread weakly has the sign of
\(1-2\omega\). The effect is strict when \(\E[(R-\ell)_+]\) changes strictly and \(\omega\ne1/2\). At \(\omega=1/2\), the
spread equals \((h-\ell)/2\) at every incumbent distribution.
\end{proposition}

\begin{proof}[Proof of Proposition~\ref{prop:payment}]
\emph{The continuation.} After preparation, the bidder with the largest value, \(\max\{R,\theta\}\), acquires the target,
and the next bidder's value, \(\min\{R,\theta\}\), is the enforceable fallback; without a rival, the fallback is
zero. For \(\max\{R,\theta\}>\min\{R,\theta\}\) and \(0<\omega<1\), Nash bargaining maximizes
\((T_\omega-\min\{R,\theta\})^\omega(\max\{R,\theta\}-T_\omega)^{1-\omega}\) over
\(\min\{R,\theta\}\le T_\omega\le\max\{R,\theta\}\). The logarithm has derivative
\(\omega/(T_\omega-\min\{R,\theta\})-(1-\omega)/(\max\{R,\theta\}-T_\omega)\) and a strictly negative second
derivative, so the unique optimum solves \(\omega(\max\{R,\theta\}-T_\omega)=(1-\omega)(T_\omega-\min\{R,\theta\})\),
that is \(T_\omega=(1-\omega)\min\{R,\theta\}+\omega\max\{R,\theta\}\). At \(R=\theta\), feasibility fixes
\(T_\omega=\theta\), and at \(\omega=0\) the next bidder accepts the fallback, so \(T_\omega=\min\{R,\theta\}\). Hence the
transfer is \(T_\omega(R,\theta)=(1-\omega)\min\{R,\theta\}+\omega\max\{R,\theta\}\).

\emph{Profit and spread.} A rival wins when \(\theta>R\) and keeps \(\theta-T_\omega=(1-\omega)(\theta-R)\), so
\(g_{\theta,\omega}=(1-\omega)\E[(\theta-R)_+]\). For \(R\le\ell\), both classes win, and
\(T_\omega(R,h)-T_\omega(R,\ell)=\omega(h-\ell)\). For \(\ell<R<h\), the high class wins and the low class loses, and
\(T_\omega(R,h)-T_\omega(R,\ell)=(1-\omega)R+\omega h-(1-\omega)\ell-\omega R=\omega(h-\ell)+(1-2\omega)(R-\ell)\). The
support lies below \(h\), so the average over the two cases gives
\(\Delta_\omega=\omega(h-\ell)+(1-2\omega)\E[(R-\ell)_+]\).

\emph{Signs.} The function \((\theta-R)_+\) is nonincreasing in \(R\), and \((R-\ell)_+\) is nondecreasing. A
first-order stochastic strengthening of the incumbent therefore weakly lowers \(g_{\theta,\omega}\) and weakly moves
\(\E[(R-\ell)_+]\) up, so the spread moves with the sign of \(1-2\omega\), strictly when \(\E[(R-\ell)_+]\) rises
strictly and \(\omega\ne1/2\). At \(\omega=1/2\), the coefficient \(1-2\omega\) is zero, and \(\Delta_\omega=(h-\ell)/2\).
\end{proof}

For \(R\sim U[0,r]\), \(\E[(R-\ell)_+]=(r-\ell)^2/(2r)=\DT(r)\), so \(\Delta_\omega(r)=\omega(h-\ell)+(1-2\omega)\DT(r)\).
In the example, \(\Delta_\omega(r_1)-\Delta_\omega(r_0)=(1-2\omega)[\DT(r_1)-\DT(r_0)]=\reg{pay_gap_coef}\,(1-2\omega)\),
which is zero exactly at \(\omega=1/2\). The proposition fixes the acquisition stage after every history.

\subsubsection*{Reserves}

At a general reserve, a class value below \(p\) gives no admissible bid. For \(R\sim U[0,r]\) and a point value \(\theta\),
integration over \(R<p\), \(p\le R<\theta\) and \(R\ge\theta\) gives the kernel
\begin{equation}\label{eq:a7-kernel}
\bigl(t(\theta),\,g(\theta),\,F(\theta)\1\{\theta\ge p\}\bigr)=
\begin{cases}
(t_0,0,0),&\theta<p,\\
(p,\,\theta-p,\,1),&r\le p\le\theta,\\
\bigl(\theta-(\theta^2-p^2)/(2r),\,(\theta^2-p^2)/(2r),\,\theta/r\bigr),&p\le\theta\le r,\\
\bigl(r/2+p^2/(2r),\,\theta-r/2-p^2/(2r),\,1\bigr),&p\le r\le\theta,
\end{cases}
\end{equation}
with \(t_0=p(1-p/r)_+\). The pieces agree at \(\theta=r\) and at \(p=r\). A class law averages the kernel. Both
\(t(\theta)\) and \(g(\theta)\) are nondecreasing in \(\theta\), and \(t(\theta)\ge t_0\); under the class ranking, \(\DT\ge0\), \(w_L\ge0\) and
\(g_H\ge g_L\ge0\). For \(\ell<p<r<h\) with point values, \(t_L=t_0=p-p^2/r\), \(g_L=0\), \(t_H=r/2+p^2/(2r)\),
\(g_H=h-t_H\) and \(\DT=r/2-p+3p^2/(2r)\). For \(r<p<h\), \(t_0=t_L=g_L=0\), \(t_H=p\), \(g_H=h-p\) and \(\DT=p\).

\begin{proposition}[Reserves]\label{prop:reserves}
Fix a reserve, and allow class values below it. Keep the class ranking, the costs, the slot and a strengthening
\(F_1\le F_0\), and compute each class object with the reserve. Suppose \(k<\rho\kappa\min\{\Delta_0,\Delta_1\}\).
\begin{enumerate}[label=(\alph*),leftmargin=2em]
\item If \(B_0(\pi)<c_H<B_1(M)\), entry is the same at the two strengths:
\(\Ent_0=\Ent_1=\rho+(1-\rho)(b+\pi\lambda)\).
\item If \(B_0(\pi)<c_H<B_0(M)\) and \(B_1(M)<c_H\), entry falls strictly:
\(\Ent_0=\rho+(1-\rho)(b+\pi\lambda)>\rho=\Ent_1\).
\end{enumerate}
Each economy has a unique outcome under every pair of tie rules.
\end{proposition}

\begin{proof}[Proof of Proposition~\ref{prop:reserves}]
\emph{Orders.} The zero-cost rival prepares after every price, so entry is at least \(\rho\). Lemma~\ref{lem:trading}
uses only (E1), (E4), (E6), \(e\in[\rho,1]\) and \(\DT\ge0\), so it holds with the reserve kernel. Every correct order earns
at least \(\rho\kappa\Delta_i-k>0\) under every candidate, and every wrong-signed order at most \(-k\). Full orders are
forced at both strengths, and the posteriors are \((m,\pi,M)\).

\emph{Entry.} Gross profit falls with incumbent strength at every belief, because \((\theta-\max\{p,R\})_+\) is
nonincreasing in \(R\). In part (a), \(B_1(\pi)\le B_0(\pi)<c_H<B_1(M)\le B_0(M)\), so at both strengths only the
order \(1\) recruits, and \(B_i(m)\le B_i(\pi)\) keeps the sale below the cost. In part (b), the weak strength has
the same pattern, and at the strong strength \(B_1(\mu)\le B_1(M)<c_H\) at every posterior. Every realized comparison
is strict.

\emph{The tuples.} At each strength the tuple has full orders, posteriors \((m,\pi,M)\), entry \((\rho,\rho,1)\) or
\((\rho,\rho,\rho)\) as above, \(\zeta(P(x))=[e(x)-\rho]/(1-\rho)\), prices
\(P_i(x)=t_{0,i}+e_i(x)\,w_i(\mu(x))\), \(\bar\mu(P_i(x))=\mu(x)\), and truthful bids. Since \(\Delta_i>0\) and entry is
nondecreasing in the posterior, the three prices differ, so (E2) holds. (E1), (E3) and (E4) hold by the computations
above, (E5) by the strict bounds, and (E6) by truthful bids. The bounds hold for every candidate, so each tuple is
unique, under every pair of tie rules. A purchase has probability \(b+\pi\lambda\), which gives both entry formulas.
\end{proof}

\paragraph{Reserves in the example.}
At each reserve and strength below, three margins decide the outcome: the trading margin \(\rho\kappa\DT-k\), the
prior-cost margin \(c_H-B(1/2)\) and the top-posterior margin \(B(M)-c_H\). Each follows from the kernel formulas
above, and the triples list them in this order.
\begin{itemize}[leftmargin=2em]
\item \emph{A reserve that excludes the low class.} At the reserve \(\reg{rsv_points_excl}\) with point values, the low
class lies below the reserve and the high class above both incumbent supports, so the formulas above apply. The
margins are \((\reg{p2_reserve_points_exclude_r0_global},\reg{p2_reserve_points_exclude_r0_prior_cost_margin},
\reg{p2_reserve_points_exclude_r0_cost_sign_plus})\) at \(r_0\) and
\((\reg{p2_reserve_points_exclude_r1_global},\reg{p2_reserve_points_exclude_r1_prior_cost_margin},
\reg{p2_reserve_points_exclude_r1_cost_sign_plus})\) at \(r_1\). Part (a) applies, and entry is
\(\reg{p2_reserve_points_exclude_r0_entry}\) at both strengths. Proceeds exceed their levels at the reserve
\(\reg{prim_p}\) by \(\reg{margin_p2_reserve_points_exclude_gain0_B}\) and \(\reg{margin_p2_reserve_points_exclude_gain1_B}\).
For the narrow bands at \(\reg{rsv_bands_excl}\), the margins are
\((\reg{p2_reserve_bands_exclude_r0_global},\reg{p2_reserve_bands_exclude_r0_prior_cost_margin},
\reg{p2_reserve_bands_exclude_r0_cost_sign_plus})\) and
\((\reg{p2_reserve_bands_exclude_r1_global},\reg{p2_reserve_bands_exclude_r1_prior_cost_margin},
\reg{p2_reserve_bands_exclude_r1_cost_sign_plus})\), with proceeds gains \(\reg{margin_p2_reserve_bands_exclude_gain0_B}\)
and \(\reg{margin_p2_reserve_bands_exclude_gain1_B}\).
\item \emph{Deterrence.} At the reserve \(\reg{rsv_deter}\), the reserve exceeds \(r_0\), so the weak incumbent never
meets it, and at \(r_0\) the objects are \(t_0=t_L=g_L=0\), \(t_H=p\) and \(g_H=h-p\). The margins are
\((\reg{p2_reserve_points_deter_r0_global},\reg{p2_reserve_points_deter_r0_prior_cost_margin},
\reg{p2_reserve_points_deter_r0_cost_sign_plus})\) at \(r_0\) and
\((\reg{p2_reserve_points_deter_r1_global},\reg{p2_reserve_points_deter_r1_prior_cost_margin},
\reg{p2_reserve_points_deter_r1_cost_sign_plus})\) at \(r_1\). Part (b) applies: entry is
\(\reg{p2_reserve_points_deter_r0_entry}\) at \(r_0\) and \(\reg{p2_reserve_points_deter_r1_entry}\) at \(r_1\). Proceeds
exceed their levels at the reserve \(\reg{prim_p}\) by \(\reg{margin_p2_reserve_points_rev0_B}\) and
\(\reg{margin_p2_reserve_points_rev1_B}\). The point values and the narrow bands give the same objects here, because the
reserve lies above the low band and below the high band. The strong upper profit equals \(c_H\) at
\(p_H=\sqrt{2r_1(h-c_H/M)-r_1^2}\approx\reg{p2_reserve_ph}\), and the weak upper profit equals \(c_H\) at
\(h-c_H/M\approx\reg{p2_reserve_weak_top}\); the reserve lies inside this interval, at distances
\(\reg{margin_p2_deter_lower_B}\) and \(\reg{margin_p2_deter_upper_B}\) from its ends, and \(p_H-r_0=\reg{margin_p2_ph_weak_B}\).
\item \emph{Why deterrence excludes the weak incumbent.} For point values and \(1<p\le r_0\), the formulas above give
\(\DT(r_0)-\reg{rsv_Delta_r0_at_ell}=(p-1)(5p+1)/4\ge0\) and \(\DT(r_1)-1=(p-1)^2/2\ge0\), so the trading margins stay at least
\(\reg{margin_p2_reserve_uniform_weak_B}\) and \(\reg{margin_p2_reserve_uniform_strong_B}\). The strong upper profit is
\(M(h-r_1/2-p^2/(2r_1))\), which at \(p=r_0\) exceeds \(c_H\) by \(\reg{margin_p2_reserve_admissible_top_B}\) and is larger
at every lower reserve. The weak upper profit is larger still, and both prior profits lie below \(c_H\). So both
strengths recruit after a purchase only, and entry is flat. For the narrow bands with
\(\reg{p2_input_bands_low_hi}<p\le r_0\), the low band lies below \(p\) and the high band above \(r_1\), so every band object equals the
point-value object, and the same conclusion holds. Strict deterrence by a reserve that excludes the low class therefore
needs \(p>r_0\).
\item \emph{Boundary reserves.} At \(p=p_H\), the purchase leaves the high-cost rival indifferent at \(r_1\):
when an indifferent high-cost rival prepares, entry is \((\rho,\rho,1)\), and when it may prepare with any
probability, the last component takes any value in \([\rho,1]\); the trading margin there is
\(\reg{margin_p2_ph_trade_B}\). At \(p=h\) with point values, every acquisition profit is zero, the definition fixes
zero-cost preparation, \(\DT=h\), the bound \(h\rho\kappa>k\) forces full orders, and entry is
\(\reg{p2_endpoint_r1_entry}\) with proceeds \(\reg{p2_endpoint_r1_proceeds}\).
\end{itemize}
